\documentclass[11pt]{article}
\usepackage[T1]{fontenc}
\usepackage[utf8]{inputenc}
\usepackage{lmodern}
\usepackage{amsmath,amssymb,amsthm}
\usepackage{booktabs}
\usepackage{enumitem}
\usepackage[margin=1in]{geometry}
\usepackage[hidelinks]{hyperref}

%======================================================================
%  Blueprint metadata macros.
%
%  This file is written so that another tool/LLM can turn it into a
%  `leanblueprint` document and drive auto-formalization.
%
%  * \uses{lab1,lab2,...}  records the dependency edges (the blueprint DAG).
%  * \leanname{Ident}      records the *intended* Lean declaration name
%                          (metadata only -- NOT Lean code). The names of the
%                          three core definitions and of the final theorem are
%                          FIXED by `Showcase_ThmA.lean` and must be reused
%                          verbatim; all other names are suggestions.
%  * \mlib{Name}           marks a fact that is already in Mathlib (a DAG leaf,
%                          i.e. `\leanok` from the start).
%
%  To build the real blueprint, replace these three definitions by the
%  blueprint package's own \uses / \lean / \leanok, or keep them: they are
%  compile-safe on their own.
%======================================================================
\newcommand{\uses}[1]{\par\smallskip{\normalfont\scriptsize$\vdash$\ \textsf{Uses:} \texttt{\detokenize{#1}}}}
\newcommand{\leanname}[1]{\ {\normalfont\scriptsize[\textsf{Lean:} \texttt{#1}]}}
\newcommand{\mlib}[1]{{\normalfont\scriptsize[\textsf{Mathlib:} \texttt{#1}]}}

\theoremstyle{definition}
\newtheorem{definition}{Definition}[section]
\newtheorem{construction}[definition]{Construction}
\newtheorem{convention}[definition]{Convention}
\theoremstyle{plain}
\newtheorem{theorem}[definition]{Theorem}
\newtheorem{proposition}[definition]{Proposition}
\newtheorem{lemma}[definition]{Lemma}
\newtheorem{corollary}[definition]{Corollary}
\newtheorem{fact}[definition]{Imported fact}
\theoremstyle{remark}
\newtheorem{remark}[definition]{Remark}

\newcommand{\N}{\mathbb{N}}
\newcommand{\R}{\mathbb{R}}
\newcommand{\norm}[1]{\lVert #1\rVert}
\newcommand{\abs}[1]{\lvert #1\rvert}
\newcommand{\pos}[1]{{#1}^{+}}
\newcommand{\Xp}{X_{+}}
\newcommand{\sotends}{\xrightarrow{\sigma o}}
\newcommand{\dzero}{\mathbin{\downarrow}0}
\newcommand{\act}{\cdot}
\newcommand{\Xrho}{X_{\rho}}
\newcommand{\bigmeet}{\bigwedge}
\newcommand{\bigjoin}{\bigvee}

\title{Theorem A ($\sigma$-order form) --- Lean-friendly blueprint edition}
\author{}
\date{}

\begin{document}
\maketitle

\begin{abstract}
This is a refactoring of the proof of Theorem A into atomic, separately
formalizable statements, annotated with an explicit dependency graph
(\texttt{\char`\\uses}), intended Lean identifiers (\texttt{[Lean: \dots]}),
and the Mathlib results each step relies on (\texttt{[Mathlib: \dots]}).
It is meant to be handed to a blueprint builder / auto-formalization agent.
The mathematics is unchanged from the source note; only the packaging is new.
No Lean code appears here.
\end{abstract}

\tableofcontents

%======================================================================
\section{How to read this document}
%======================================================================

Each numbered block is one intended Lean declaration. A block's
\texttt{\char`\\uses} line lists the labels of the blocks it depends on; together
these lines are the blueprint DAG. Blocks marked \emph{Imported fact} are
already available in Mathlib (verified against Mathlib commit
\texttt{d90090f647}, Lean \texttt{v4.31.0-rc2}) and should be treated as leaves.

\paragraph{Two things that are not off-the-shelf.}
\begin{itemize}[leftmargin=1.5em]
\item Lemma~\ref{lem:abs_smul} (\(\abs{\alpha \act x}=\abs{\alpha}\act\abs{x}\) on a
      \emph{non-linearly-ordered} vector lattice) is the \emph{only} genuinely
      missing Mathlib lemma. A complete elementary proof is given
      (Lemma~\ref{lem:smul_sup_nonneg} + Lemma~\ref{lem:abs_smul}).
\item The auxiliary norm \(\rho\) must live on a \emph{second} normed structure
      over the same vector space (Construction~\ref{con:Xrho}); Lean cannot carry
      two \texttt{NormedAddCommGroup} instances on one type, so a type synonym is
      unavoidable. This is standard packaging, not new theory.
\end{itemize}

\begin{convention}[Indexing and notation]\label{conv:index}
All sequences are indexed by \(\N=\{0,1,2,\dots\}\). For a scalar sequence
\(a:\N\to\R\) the \(m\)-th partial sum is
\[
  s_m(a)\;=\;\sum_{k=0}^{m} a_k e_k
  \qquad\text{(Lean: \(\textstyle\sum_{k\in\mathrm{range}(m+1)} a_k\act e_k\)).}
\]
``\(\forall^{\infty} n\)'' means ``for all sufficiently large \(n\)''
(Lean: \(\exists n_0,\ \forall n\ge n_0\)). \(\Xp=\{x\in X: 0\le x\}\).
The absolute value is \(\abs{x}=x\vee(-x)\) and the positive part is
\(\pos{x}=x\vee 0\). Scalar action of \(\alpha\in\R\) on \(x\in X\) is written
\(\alpha\act x\) (Lean: \(\alpha\act x\)).
\end{convention}

\begin{remark}[The ambient structure]\label{rem:structure}
Throughout, \(X\) is a Banach lattice, encoded exactly as in
\texttt{Showcase\_ThmA.lean} by the instance bundle
\[
  \texttt{NormedAddCommGroup } X,\ \texttt{Lattice } X,\
  \texttt{HasSolidNorm } X,\ \texttt{IsOrderedAddMonoid } X,
\]
\[
  \texttt{NormedSpace } \R\, X,\ \texttt{PosSMulMono } \R\, X,\
  \texttt{CompleteSpace } X .
\]
This is precisely Mathlib's normed-lattice setting together with a real module
whose positive scalars act monotonically, plus completeness. All ``Imported
fact'' blocks below hold for this exact bundle.
\end{remark}

%======================================================================
\section{Imported facts from Mathlib (DAG leaves)}\label{sec:mathlib}
%======================================================================

These require no proof; they fix the names the formalization should call.

\begin{fact}[Solid norm is monotone]\label{fact:solid}
\mlib{norm\_le\_norm\_of\_abs\_le\_abs, HasSolidNorm.solid}
For \(a,b\in X\), \(\abs{a}\le\abs{b}\Rightarrow \norm{a}\le\norm{b}\); also
\(\norm{\,\abs{a}\,}=\norm{a}\) (\mlib{norm\_abs\_eq\_norm}). In particular
\(0\le a\le b\Rightarrow \norm{a}\le\norm{b}\).
\end{fact}

\begin{fact}[Order topology is closed]\label{fact:orderclosed}
\mlib{HasSolidNorm.orderClosedTopology}
\(X\) has \texttt{OrderClosedTopology}: for \(a\in X\) the sets
\(\{y:y\le a\}\), \(\{y:a\le y\}\) and \([-a,a]\) are norm-closed. Consequently,
if \(y_i\to y\) in norm and \(y_i\le a\) for all \(i\), then \(y\le a\)
(\mlib{le\_of\_tendsto', ge\_of\_tendsto', isClosed\_Icc}).
\end{fact}

\begin{fact}[Continuity of the lattice maps]\label{fact:contabs}
\mlib{continuous\_posPart, continuous\_negPart, HasSolidNorm.continuousSup}
The maps \(x\mapsto\pos{x}\), \(x\mapsto x^-\), \(x\mapsto\abs{x}\),
\(\vee\), \(\wedge\) are norm-continuous.
\end{fact}

\begin{fact}[Lattice-ordered group algebra]\label{fact:absgroup}
\mlib{abs\_add\_le, le\_abs\_self, neg\_abs\_le, abs\_le, abs\_nonneg, abs\_neg, posPart\_nonneg, posPart\_eq\_zero}
For \(a,b\in X\): \(\abs{a+b}\le\abs{a}+\abs{b}\); \(a\le\abs{a}\);
\(\abs{a}\le b \Leftrightarrow -b\le a\le b\); \(0\le\abs{a}\); \(\abs{-a}=\abs{a}\);
\(0\le\pos{a}\); and \(\pos{a}=0 \Leftrightarrow a\le 0\).
(All valid for the non-linear lattice-ordered group \(X\).)
\end{fact}

\begin{fact}[Monotone positive scalar action]\label{fact:possmul}
\mlib{PosSMulMono, monotone\_smul\_left\_of\_nonneg, smul\_nonneg}
For \(0\le\alpha\) and \(a\le b\): \(\alpha\act a\le\alpha\act b\); and
\(0\le\alpha,\ 0\le a\Rightarrow 0\le\alpha\act a\).
\end{fact}

\begin{fact}[Absolutely convergent series and tails]\label{fact:series}
\mlib{Summable.of\_norm, tsum\_nonneg, tsum\_eq\_zero\_add, tendsto\_sum\_nat\_add}
If \(\sum_r\norm{v_r}<\infty\) then \(\sum_r v_r\) converges in the Banach space
\(X\); if moreover \(v_r\ge 0\) then \(\sum_{r\ge R} v_r\ge 0\),
\(\sum_{r\ge R}v_r = v_R+\sum_{r\ge R+1}v_r\), and
\(\norm{\sum_{r\ge R}v_r}\to 0\) as \(R\to\infty\).
\end{fact}

\begin{fact}[Real infima]\label{fact:csinf}
\mlib{csInf\_le, le\_csInf, Real.sInf\_nonneg}
For a nonempty \(S\subseteq\R\) bounded below: \(\inf S\le s\) for \(s\in S\),
and \(c\le\inf S\) iff \(c\) is a lower bound.
\end{fact}

\begin{fact}[Finite lattice sup/inf are least/greatest bounds]\label{fact:finset}
\mlib{Finset.le\_sup', Finset.sup'\_le, Finset.le\_inf', Finset.inf'\_le, Finset.inf'\_mono}
For a nonempty finite family, \(\bigjoin\) is its least upper bound and
\(\bigmeet\) its greatest lower bound, both monotone in the family.
\end{fact}

\begin{fact}[Finite-dimensional subspaces]\label{fact:findim}
\mlib{Submodule.closed\_of\_finiteDimensional, LinearMap.continuous\_of\_finiteDimensional, Basis.mk, LinearIndependent}
A finite-dimensional subspace of \(X\) is norm-closed; every linear map out of a
finite-dimensional normed space is continuous; a finite linearly independent
family is a basis of its span, giving continuous coordinate maps.
\end{fact}

\begin{fact}[Norm from a norm-core; bounded inverse theorem]\label{fact:banach}
\mlib{NormedAddCommGroup.ofCore, AddGroupNorm.toNormedAddCommGroup, LinearMap.mkContinuous, ContinuousLinearEquiv.ofBijective, ContinuousLinearMap.le\_opNorm}
A function satisfying the norm axioms induces a \texttt{NormedAddCommGroup}. If
\(E\) is a Banach space and \(f:E\to F\) is a continuous linear bijection onto a
normed space \(F\), then \(f^{-1}\) is continuous (\texttt{ofBijective} with
\(\ker f=\bot\), \(\operatorname{range}f=\top\)); \(\norm{f^{-1}y}\le\norm{f^{-1}}\,\norm{y}\).
\end{fact}

\begin{fact}[Completeness via absolute convergence]\label{fact:completecrit}
\mlib{completeSpace criterion: ``summable-in-norm $\Rightarrow$ has-sum'' (e.g.\ via \texttt{Metric.complete\_of\_cauchySeq\_tendsto} or the normed-group series criterion)}
A normed space is complete once every norm-summable series converges. Used in
Proposition~\ref{prop:rho_complete}; an equivalent self-contained subsequence
argument is described there.
\end{fact}

%======================================================================
\section{Preliminaries I: scalars versus the lattice}\label{sec:scalars}
%======================================================================

This section supplies the one lemma missing from Mathlib and its helper.

\begin{lemma}[Positive scalars distribute over joins]\label{lem:smul_sup_nonneg}
\leanname{smul\_sup\_of\_nonneg}
For \(\alpha\in\R\) with \(0\le\alpha\) and \(a,b\in X\),
\(\alpha\act(a\vee b)=(\alpha\act a)\vee(\alpha\act b)\).
\uses{fact:possmul}
\end{lemma}
\begin{proof}
Write \((\ast_d)\) for the one-sided inequality
\((d\act p)\vee(d\act q)\le d\act(p\vee q)\), valid for \(0\le d\): indeed
\(p\le p\vee q\) gives \(d\act p\le d\act(p\vee q)\) by Fact~\ref{fact:possmul},
similarly for \(q\), and take the join.

Applying \((\ast_\alpha)\) gives ``\(\ge\)'' of the claim (i.e.\ \((\alpha\act a)\vee(\alpha\act b)\le\alpha\act(a\vee b)\)).
For ``\(\le\)'': if \(\alpha=0\) both sides are \(0\). If \(\alpha>0\), apply
\((\ast_{\alpha^{-1}})\) to \(p=\alpha\act a,\ q=\alpha\act b\):
\((\alpha^{-1}\act(\alpha\act a))\vee(\alpha^{-1}\act(\alpha\act b))\le
\alpha^{-1}\act\big((\alpha\act a)\vee(\alpha\act b)\big)\),
whose left side is \(a\vee b\). Act by \(\alpha\ge 0\) (monotone) and use
\(\alpha\act(\alpha^{-1}\act y)=y\):
\(\alpha\act(a\vee b)\le(\alpha\act a)\vee(\alpha\act b)\).
\end{proof}

\begin{lemma}[Absolute value is scalar-homogeneous]\label{lem:abs_smul}
\leanname{abs\_smul\_vectorLattice}
For all \(\alpha\in\R\) and \(x\in X\), \(\abs{\alpha\act x}=\abs{\alpha}\act\abs{x}\).
\uses{lem:smul_sup_nonneg,fact:absgroup}
\end{lemma}
\begin{proof}
If \(\alpha\ge 0\): using \(-(\alpha\act x)=\alpha\act(-x)\) and
Lemma~\ref{lem:smul_sup_nonneg},
\[
  \abs{\alpha\act x}=(\alpha\act x)\vee(\alpha\act(-x))
  =\alpha\act\big(x\vee(-x)\big)=\alpha\act\abs{x}=\abs{\alpha}\act\abs{x}.
\]
If \(\alpha\le 0\): put \(d=-\alpha\ge 0\); then \(\alpha\act x=-(d\act x)\), so
\(\abs{\alpha\act x}=\abs{d\act x}=d\act\abs{x}=\abs{\alpha}\act\abs{x}\) using
\(\abs{-y}=\abs{y}\) and the previous case.
\end{proof}

\begin{remark}\label{rem:gap}
Mathlib's \texttt{abs\_smul} and \texttt{smul\_max\_of\_nonneg} require
\texttt{LinearOrder} on the module, which fails for a general Banach lattice; for
general lattices Mathlib supplies only the one-sided \texttt{smul\_inf\_le}. Hence
Lemmas~\ref{lem:smul_sup_nonneg}--\ref{lem:abs_smul} are the sole pieces of new
theory required.
\end{remark}

%======================================================================
\section{Preliminaries II: regulators (the calculus of \texorpdfstring{$z\dzero$}{z down 0})}\label{sec:reg}
%======================================================================

\begin{definition}[Regulator, \(z\dzero\)]\label{def:downto}
\leanname{DownToZero}
A sequence \(z:\N\to X\) is a \emph{regulator} (written \(z\dzero\)) if
(i) \(0\le z_m\) for all \(m\); (ii) \(z_{m+1}\le z_m\) for all \(m\); and
(iii) whenever \(0\le h\) and \(h\le z_m\) for all \(m\), then \(h=0\).
\end{definition}

\begin{lemma}[Zero is a regulator]\label{lem:reg_zero}
\leanname{downToZero\_zero}
The constant sequence \(0\) satisfies \(0\dzero\).
\uses{def:downto}
\end{lemma}
\begin{proof}
(i),(ii) are trivial. For (iii): \(0\le h\le 0\) forces \(h=0\) by antisymmetry.
\end{proof}

\begin{lemma}[Regulators are closed under addition]\label{lem:reg_add}
\leanname{DownToZero.add}
If \(z\dzero\) and \(w\dzero\) then \((z_m+w_m)_m\dzero\).
\uses{def:downto,fact:absgroup}
\end{lemma}
\begin{proof}
Positivity and monotonicity are termwise (add nonnegatives; add the two
decrease inequalities). For (iii), suppose \(0\le h\le z_m+w_m\) for all \(m\).
Fix \(M\). For every \(m\ge M\), monotonicity of \(z\) gives
\(h\le z_M+w_m\), i.e.\ \(h-z_M\le w_m\); since also \(0\le w_m\) we get
\(\pos{(h-z_M)}\le w_m\). For \(m<M\), \(w_m\ge w_M\ge \pos{(h-z_M)}\), so in fact
\(\pos{(h-z_M)}\le w_m\) for \emph{all} \(m\). As \(w\dzero\) and
\(\pos{(h-z_M)}\ge 0\), we get \(\pos{(h-z_M)}=0\), i.e.\ \(h\le z_M\)
(Fact~\ref{fact:absgroup}, \texttt{posPart\_eq\_zero}). This holds for every \(M\)
and \(h\ge0\), so \(h=0\) since \(z\dzero\).
\end{proof}

\begin{lemma}[Regulators are closed under scaling]\label{lem:reg_smul}
\leanname{DownToZero.smul}
If \(z\dzero\) and \(\alpha\in\R\), then \((\abs{\alpha}\act z_m)_m\dzero\).
\uses{def:downto,lem:reg_zero,fact:possmul}
\end{lemma}
\begin{proof}
Positivity and monotonicity follow from Fact~\ref{fact:possmul}
(\(\abs{\alpha}\ge0\)). If \(\alpha=0\) the sequence is \(0\), covered by
Lemma~\ref{lem:reg_zero}. If \(\alpha\neq0\), let \(0\le h\le\abs{\alpha}\act z_m\)
for all \(m\). Acting by \(\abs{\alpha}^{-1}\ge0\) gives
\(0\le\abs{\alpha}^{-1}\act h\le z_m\) for all \(m\); since \(z\dzero\),
\(\abs{\alpha}^{-1}\act h=0\), hence \(h=\abs{\alpha}\act(\abs{\alpha}^{-1}\act h)=0\).
\end{proof}

\begin{lemma}[Diagonal regulator]\label{lem:diagonal}
\leanname{sigmaOrderTendsto\_of\_diagonal}
Let \(y:\N\to X\) and \(x\in X\). Suppose there are
\begin{enumerate}[label=(\alph*),leftmargin=2.2em]
\item \(c:\N\to X\) with \(c\dzero\); and
\item for each \(R\), a regulator \(w^{R}\dzero\);
\end{enumerate}
such that for all \(R,j\) one has \(\abs{y_m-x}\le c_R+w^{R}_{j}\) for all
sufficiently large \(m\). Then \(y_m\sotends x\) (Definition~\ref{def:sigo}).
\uses{def:downto,def:sigo,lem:reg_add,fact:absgroup,fact:finset}
\end{lemma}
\begin{proof}
Define \(b_j=\bigmeet_{R=0}^{j}\big(c_R+w^{R}_{j}\big)\).

\emph{Positivity/monotonicity.} Each \(c_R+w^R_j\ge0\), so \(b_j\ge0\)
(Fact~\ref{fact:finset}). For monotonicity, since \(w^R_{j+1}\le w^R_j\), each
term of \(b_{j+1}\) with \(R\le j\) is \(\le\) the corresponding term of \(b_j\),
and \(b_{j+1}\) has one extra term; hence \(b_{j+1}\le b_j\)
(Fact~\ref{fact:finset}).

\emph{The bound.} Fix \(j\). By hypothesis, for each \(R\in\{0,\dots,j\}\) there is
\(m_0(R)\) beyond which \(\abs{y_m-x}\le c_R+w^R_j\); for \(m\ge\max_R m_0(R)\) all
these hold, so \(\abs{y_m-x}\le b_j\) (greatest lower bound, Fact~\ref{fact:finset}).

\emph{\(b\dzero\).} Suppose \(0\le h\le b_j\) for all \(j\). Fix \(R\). For \(j\ge R\),
\(b_j\le c_R+w^R_j\), so \(h\le c_R+w^R_j\), giving \(\pos{(h-c_R)}\le w^R_j\) for
\(j\ge R\), hence (monotonicity of \(w^R\)) for all \(j\). As \(w^R\dzero\),
\(\pos{(h-c_R)}=0\), i.e.\ \(h\le c_R\); true for all \(R\), so \(h=0\) since
\(c\dzero\). Thus \(b\dzero\) and it regulates \(y_m\sotends x\).
\end{proof}

%======================================================================
\section{\texorpdfstring{$\sigma$}{sigma}-order convergence (basic calculus)}\label{sec:so}
%======================================================================

\begin{definition}[\(\sigma\)-order convergence]\label{def:sigo}
\leanname{SigmaOrderTendsto}
\(x:\N\to X\) \emph{\(\sigma\)-order converges} to \(x_0\), written \(x_n\sotends x_0\),
if there is a regulator \(z\dzero\) with
\(\forall m,\ \forall^{\infty} n,\ \abs{x_n-x_0}\le z_m\).
\uses{def:downto}
\end{definition}

\begin{lemma}[Eventually constant \(\Rightarrow\) converges]\label{lem:so_eventuallyEq}
\leanname{sigmaOrderTendsto\_of\_eventuallyEq}
If \(\forall^{\infty} n,\ x_n=x_0\), then \(x_n\sotends x_0\).
\uses{def:sigo,lem:reg_zero}
\end{lemma}
\begin{proof}
Use the regulator \(0\) (Lemma~\ref{lem:reg_zero}); for each \(m\), eventually
\(\abs{x_n-x_0}=0\le 0\).
\end{proof}

\begin{lemma}[Additivity]\label{lem:so_add}
\leanname{SigmaOrderTendsto.add}
If \(x_n\sotends x\) and \(y_n\sotends y\), then \(x_n+y_n\sotends x+y\).
\uses{def:sigo,lem:reg_add,fact:absgroup}
\end{lemma}
\begin{proof}
Let \(z,w\) be regulators for \(x,y\). By Lemma~\ref{lem:reg_add}, \(z+w\dzero\).
For each \(m\), on the eventual set where both
\(\abs{x_n-x}\le z_m\) and \(\abs{y_n-y}\le w_m\) hold,
\(\abs{(x_n+y_n)-(x+y)}\le\abs{x_n-x}+\abs{y_n-y}\le z_m+w_m\).
\end{proof}

\begin{lemma}[Homogeneity]\label{lem:so_smul}
\leanname{SigmaOrderTendsto.smul}
If \(x_n\sotends x\) and \(\alpha\in\R\), then \(\alpha\act x_n\sotends \alpha\act x\).
\uses{def:sigo,lem:reg_smul,lem:abs_smul,fact:possmul}
\end{lemma}
\begin{proof}
Let \(z\dzero\) regulate \(x\). By Lemma~\ref{lem:reg_smul},
\((\abs{\alpha}\act z_m)_m\dzero\). For each \(m\), eventually
\(\abs{\alpha\act x_n-\alpha\act x}=\abs{\alpha\act(x_n-x)}
=\abs{\alpha}\act\abs{x_n-x}\le\abs{\alpha}\act z_m\)
by Lemma~\ref{lem:abs_smul} and Fact~\ref{fact:possmul}.
\end{proof}

\begin{lemma}[\(\sigma\)-order limits are order bounded]\label{lem:so_bounded}
\leanname{SigmaOrderTendsto.orderBounded}
If \(x_n\sotends x\), then there is \(u\in\Xp\) with \(\abs{x_n}\le u\) for all \(n\).
\uses{def:sigo,fact:absgroup,fact:finset}
\end{lemma}
\begin{proof}
Take a regulator \(z\). For \(m=0\) there is \(n_0\) with \(\abs{x_n-x}\le z_0\)
for \(n\ge n_0\), whence \(\abs{x_n}\le\abs{x}+z_0\) there. Set
\(u=\big(\bigjoin_{n<n_0}\abs{x_n}\big)\vee(\abs{x}+z_0)\); then \(u\ge0\) and
\(\abs{x_n}\le u\) for every \(n\) (Fact~\ref{fact:finset}).
\end{proof}

\begin{lemma}[Order bounds pass to \(\sigma\)-order limits]\label{lem:so_abs_le}
\leanname{SigmaOrderTendsto.abs\_le}
If \(x_n\sotends x\), \(u\in\Xp\), and \(\abs{x_n}\le u\) for all \(n\), then
\(\abs{x}\le u\).
\uses{def:sigo,fact:absgroup}
\end{lemma}
\begin{proof}
Take a regulator \(z\). For each \(m\), pick \(n\) large with
\(\abs{x_n-x}\le z_m\); then
\(\abs{x}\le\abs{x_n}+\abs{x_n-x}\le u+z_m\), so \(\pos{(\abs{x}-u)}\le z_m\).
Valid for all \(m\), so \(\pos{(\abs{x}-u)}=0\) (\(z\dzero\)), i.e.\ \(\abs{x}\le u\).
\end{proof}

\begin{lemma}[Order bounds pass to \emph{norm} limits]\label{lem:norm_limit_abs_le}
\leanname{abs\_le\_of\_tendsto}
If \(x_i\to x\) in norm and \(\abs{x_i}\le u\) for all \(i\), then \(\abs{x}\le u\).
\uses{fact:orderclosed,fact:absgroup}
\end{lemma}
\begin{proof}
\(\abs{x_i}\le u\) means \(-u\le x_i\le u\) (Fact~\ref{fact:absgroup}). Both
half-spaces are closed (Fact~\ref{fact:orderclosed}); passing to the norm limit
gives \(-u\le x\le u\), i.e.\ \(\abs{x}\le u\).
\end{proof}

\begin{lemma}[Tails of a positive summable series form a regulator]\label{lem:reg_tail}
\leanname{downToZero\_tsum\_tail}
Let \(u:\N\to X\) with \(u_r\ge0\) and \(\sum_r\norm{u_r}<\infty\), and set
\(U_R=\sum_{r\ge R} u_r\). Then \((U_R)_R\dzero\).
\uses{def:downto,fact:series,fact:solid}
\end{lemma}
\begin{proof}
The series converges (Fact~\ref{fact:series}). Then \(U_R\ge0\)
(\(\mathtt{tsum\_nonneg}\)), \(U_R=u_R+U_{R+1}\ge U_{R+1}\) (so decreasing), and
\(\norm{U_R}\to0\) (tail of a norm-convergent series). If \(0\le h\le U_R\) for
all \(R\), then \(\norm{h}\le\norm{U_R}\to0\) (Fact~\ref{fact:solid}), so
\(\norm{h}=0\) and \(h=0\).
\end{proof}

%======================================================================
\section{Coordinates and projections}\label{sec:coord}
%======================================================================

\begin{definition}[\(\sigma\)-order basis]\label{def:basis}
\leanname{IsSigmaOrderBasis}
\((e_n)\) is a \emph{\(\sigma\)-order basis} of \(X\) if for every \(x\in X\)
there is a unique \(a:\N\to\R\) with \(s_m(a)\sotends x\).
\uses{def:sigo}
\end{definition}

\emph{Standing assumption for \S\S\ref{sec:coord}--\ref{sec:max}}: \((e_n)\) is a
fixed \(\sigma\)-order basis of \(X\).

\begin{construction}[Coordinates and projections]\label{con:coord}
\leanname{coord, proj}
For \(x\in X\) let \(a^{x}\) be the unique sequence from Definition~\ref{def:basis}.
Define the coordinate functionals and projections
\[
  e_k^{\#}(x)=a^{x}_k,\qquad
  P_m x=s_m(a^{x})=\sum_{k=0}^{m} e_k^{\#}(x)\,e_k .
\]
\uses{def:basis}
\end{construction}

\begin{lemma}[Projections \(\sigma\)-order converge]\label{lem:proj_sotends}
\leanname{proj\_sigmaOrderTendsto}
For every \(x\), \(P_m x\sotends x\) (as \(m\to\infty\)).
\uses{con:coord,def:basis}
\end{lemma}
\begin{proof}
Immediate from Construction~\ref{con:coord}: \(P_m x=s_m(a^x)\) and
\(s_m(a^x)\sotends x\) by choice of \(a^x\).
\end{proof}

\begin{lemma}[Uniqueness of expansion]\label{lem:coord_unique}
\leanname{coord\_unique}
If \(s_m(a)\sotends x\) then \(a=a^{x}\) (equivalently \(a_k=e_k^{\#}(x)\) for all \(k\)).
\uses{con:coord,def:basis}
\end{lemma}
\begin{proof}
The uniqueness clause of Definition~\ref{def:basis}.
\end{proof}

\begin{lemma}[Coordinates are additive]\label{lem:coord_add}
\leanname{coord\_add}
\(e_k^{\#}(x+y)=e_k^{\#}(x)+e_k^{\#}(y)\) for all \(k\).
\uses{lem:proj_sotends,lem:so_add,lem:coord_unique}
\end{lemma}
\begin{proof}
\(s_m(a^{x}+a^{y})=s_m(a^{x})+s_m(a^{y})\) (linearity of the finite sum in the
coefficients) \(\sotends x+y\) by Lemma~\ref{lem:so_add}. By
Lemma~\ref{lem:coord_unique}, \(a^{x+y}=a^{x}+a^{y}\).
\end{proof}

\begin{lemma}[Coordinates are homogeneous]\label{lem:coord_smul}
\leanname{coord\_smul}
\(e_k^{\#}(\alpha\act x)=\alpha\, e_k^{\#}(x)\) for all \(k\).
\uses{lem:proj_sotends,lem:so_smul,lem:coord_unique}
\end{lemma}
\begin{proof}
\(s_m(\alpha\, a^{x})=\alpha\act s_m(a^{x})\sotends\alpha\act x\) by
Lemma~\ref{lem:so_smul}; conclude with Lemma~\ref{lem:coord_unique}.
\end{proof}

\begin{lemma}[Projections are linear]\label{lem:proj_linear}
\leanname{proj\_add, proj\_smul}
\(P_m(x+y)=P_m x+P_m y\) and \(P_m(\alpha\act x)=\alpha\act P_m x\).
\uses{lem:coord_add,lem:coord_smul}
\end{lemma}
\begin{proof}
Substitute Lemmas~\ref{lem:coord_add}--\ref{lem:coord_smul} into
\(P_m=\sum_{k=0}^m e_k^{\#}(\cdot)\,e_k\).
\end{proof}

\begin{lemma}[The basis vectors are linearly independent]\label{lem:e_indep}
\leanname{e\_linearIndependent}
The family \((e_n)\) is linearly independent.
\uses{con:coord,lem:so_eventuallyEq,lem:coord_unique}
\end{lemma}
\begin{proof}
Suppose \(\sum_{k=0}^{N} a_k e_k=0\). Let \(\hat a\) extend \((a_0,\dots,a_N)\)
by zeros. Then \(s_m(\hat a)=0\) for all \(m\ge N\), so
\(s_m(\hat a)\sotends 0\) (Lemma~\ref{lem:so_eventuallyEq}). Also the zero
sequence gives \(s_m(0)=0\sotends 0\). By uniqueness
(Lemma~\ref{lem:coord_unique}) \(\hat a=0\), i.e.\ every \(a_k=0\).
\end{proof}

%======================================================================
\section{The auxiliary norm \texorpdfstring{$\rho$}{rho}}\label{sec:rho}
%======================================================================

\begin{construction}[The functional \(\rho\)]\label{con:rho}
\leanname{rho}
For \(x\in X\) set
\[
  A(x)=\{\,u\in\Xp : \abs{P_m x}\le u \text{ for every } m\in\N\,\},
  \qquad
  \rho(x)=\inf\{\norm{u}: u\in A(x)\}.
\]
\uses{con:coord}
\end{construction}

\begin{lemma}[\(A(x)\) is nonempty and \(\rho\) is well defined]\label{lem:rho_wd}
\leanname{rho\_admissible\_nonempty}
\(A(x)\neq\varnothing\); hence \(\rho(x)\) is a well-defined nonnegative real.
\uses{con:rho,lem:proj_sotends,lem:so_bounded,fact:csinf}
\end{lemma}
\begin{proof}
\(P_m x\sotends x\) (Lemma~\ref{lem:proj_sotends}), so
Lemma~\ref{lem:so_bounded} yields \(u\ge0\) with \(\abs{P_m x}\le u\) for all
\(m\), i.e.\ \(u\in A(x)\). The set \(\{\norm{u}:u\in A(x)\}\) is nonempty and
bounded below by \(0\), so \(\rho(x)\ge0\) exists (Fact~\ref{fact:csinf}).
\end{proof}

\begin{lemma}[\(\rho\) dominates the norm]\label{lem:norm_le_rho}
\leanname{norm\_le\_rho}
\(\norm{x}\le\rho(x)\) for all \(x\).
\uses{con:rho,lem:proj_sotends,lem:so_abs_le,fact:solid,fact:csinf}
\end{lemma}
\begin{proof}
For \(u\in A(x)\): since \(P_m x\sotends x\) and \(\abs{P_m x}\le u\),
Lemma~\ref{lem:so_abs_le} gives \(\abs{x}\le u\), so \(\norm{x}\le\norm{u}\)
(Fact~\ref{fact:solid}). Thus \(\norm{x}\) is a lower bound of
\(\{\norm{u}:u\in A(x)\}\), whence \(\norm{x}\le\rho(x)\) (Fact~\ref{fact:csinf}).
\end{proof}

\begin{lemma}[\(\varepsilon\)-approximation of \(\rho\)]\label{lem:rho_approx}
\leanname{rho\_approx}
For \(x\in X\) and \(\varepsilon>0\), there is \(u\in A(x)\) with
\(\norm{u}\le\rho(x)+\varepsilon\).
\uses{con:rho,lem:rho_wd,fact:csinf}
\end{lemma}
\begin{proof}
Definition of infimum (Fact~\ref{fact:csinf}) on the nonempty set of
Lemma~\ref{lem:rho_wd}.
\end{proof}

\begin{lemma}[Homogeneity of \(\rho\)]\label{lem:rho_smul}
\leanname{rho\_smul}
\(\rho(\alpha\act x)=\abs{\alpha}\,\rho(x)\).
\uses{con:rho,lem:proj_linear,lem:abs_smul,fact:possmul}
\end{lemma}
\begin{proof}
If \(\alpha=0\), \(A(0)\ni0\) and \(\rho(0)=0\). If \(\alpha\neq0\): using
\(P_m(\alpha\act x)=\alpha\act P_m x\) (Lemma~\ref{lem:proj_linear}) and
\(\abs{\alpha\act P_m x}=\abs{\alpha}\act\abs{P_m x}\)
(Lemma~\ref{lem:abs_smul}), the map \(u\mapsto\abs{\alpha}\act u\) is a bijection
\(A(x)\to A(\alpha\act x)\) with \(\norm{\abs{\alpha}\act u}=\abs{\alpha}\norm{u}\).
Taking infima gives \(\rho(\alpha\act x)=\abs{\alpha}\rho(x)\).
\end{proof}

\begin{lemma}[Subadditivity of \(\rho\)]\label{lem:rho_add}
\leanname{rho\_add\_le}
\(\rho(x+y)\le\rho(x)+\rho(y)\).
\uses{con:rho,lem:proj_linear,lem:rho_approx,fact:absgroup,fact:csinf}
\end{lemma}
\begin{proof}
Given \(\varepsilon>0\), pick \(u\in A(x)\), \(v\in A(y)\) with
\(\norm{u}\le\rho(x)+\varepsilon\), \(\norm{v}\le\rho(y)+\varepsilon\)
(Lemma~\ref{lem:rho_approx}). Then \(u+v\ge0\) and, by
Lemma~\ref{lem:proj_linear} and Fact~\ref{fact:absgroup},
\(\abs{P_m(x+y)}=\abs{P_m x+P_m y}\le\abs{P_m x}+\abs{P_m y}\le u+v\),
so \(u+v\in A(x+y)\). Hence
\(\rho(x+y)\le\norm{u+v}\le\norm{u}+\norm{v}\le\rho(x)+\rho(y)+2\varepsilon\).
Let \(\varepsilon\to0\).
\end{proof}

\begin{lemma}[\(\rho\) is positive definite]\label{lem:rho_zero}
\leanname{rho\_eq\_zero\_iff}
\(\rho(x)=0\) iff \(x=0\).
\uses{con:rho,lem:norm_le_rho}
\end{lemma}
\begin{proof}
\(\rho(0)=0\) since \(0\in A(0)\). Conversely \(\norm{x}\le\rho(x)=0\)
(Lemma~\ref{lem:norm_le_rho}) forces \(x=0\).
\end{proof}

\begin{construction}[The normed space \((X,\rho)\)]\label{con:Xrho}
\leanname{Xrho (type synonym), instances NormedAddCommGroup / NormedSpace}
By Lemmas~\ref{lem:rho_smul}, \ref{lem:rho_add}, \ref{lem:rho_zero}, \(\rho\) is a
norm. Let \(\Xrho\) denote the vector space \(X\) equipped with \(\rho\) (a type
synonym carrying the same additive and \(\R\)-module structure). It is a normed
\(\R\)-space, and the set-theoretic identity \(\iota:\Xrho\to X\) is linear.
\uses{lem:rho_smul,lem:rho_add,lem:rho_zero,fact:banach}
\end{construction}

%======================================================================
\section{Completeness of \texorpdfstring{$(X,\rho)$}{(X,rho)}}\label{sec:complete}
%======================================================================

The heart of the argument. We first isolate the finite-dimensional step, then
prove that norm-summable series converge in \(\rho\), then conclude completeness.

\begin{lemma}[A compatible tower of finite sections is the projection tower]\label{lem:tower_eq_proj}
\leanname{proj\_of\_compatible\_tower}
Let \(x\in X\) and \(y:\N\to X\) satisfy
\begin{enumerate}[label=(\roman*),leftmargin=2.2em]
\item \(y_m\in E_m:=\operatorname{span}\{e_0,\dots,e_m\}\) for all \(m\);
\item \(Q_m y_{m+1}=y_m\), where \(Q_m:E_{m+1}\to E_m\) is the truncation
      \(\sum_{k=0}^{m+1} c_k e_k\mapsto\sum_{k=0}^{m} c_k e_k\); and
\item \(y_m\sotends x\).
\end{enumerate}
Then \(y_m=P_m x\) for all \(m\).
\uses{lem:e_indep,lem:coord_unique,fact:findim}
\end{lemma}
\begin{proof}
By linear independence (Lemma~\ref{lem:e_indep}), each \(y_m\) has unique
coordinates \(y_m=\sum_{k=0}^m c^{(m)}_k e_k\), and \(Q_m\) is well defined.
Condition (ii) gives \(c^{(m+1)}_k=c^{(m)}_k\) for \(k\le m\); set
\(a_k:=c^{(k)}_k\), so \(y_m=s_m(a)\). By (iii) and
Lemma~\ref{lem:coord_unique}, \(a=a^{x}\), i.e.\ \(y_m=P_m x\).
\end{proof}

\begin{lemma}[Norm-summable series converge in \(\rho\)]\label{lem:rho_series}
\leanname{rho\_summable\_of\_rho\_summable}
Let \((v_r)\subseteq X\) satisfy \(\sum_r\rho(v_r)<\infty\). Then the partial sums
\(x_R=\sum_{r<R} v_r\) form a \(\rho\)-Cauchy sequence that \(\rho\)-converges;
concretely there is \(x\) with \(\rho(x-x_R)\to0\).
\uses{con:rho,lem:rho_approx,lem:norm_le_rho,lem:reg_tail,lem:reg_add,lem:proj_sotends,lem:proj_linear,lem:norm_limit_abs_le,lem:diagonal,lem:tower_eq_proj,fact:series,fact:findim}
\end{lemma}
\begin{proof}
\emph{Step 1 (choose order bounds).} For each \(r\), Lemma~\ref{lem:rho_approx}
gives \(u_r\in A(v_r)\) with \(\norm{u_r}\le\rho(v_r)+2^{-r}\); thus
\(\sum_r\norm{u_r}<\infty\), and by Lemma~\ref{lem:norm_le_rho}
\(\sum_r\norm{v_r}<\infty\). Put
\[
  U_R=\sum_{r\ge R}u_r,\qquad T_R=\sum_{r\ge R}\abs{v_r}.
\]
By Lemma~\ref{lem:reg_tail}, \(U_R\dzero\) and \(T_R\dzero\) (note \(u_r,\abs{v_r}\ge0\)).

\emph{Step 2 (norm limit \(x\)).} Since \(\sum_r\norm{v_r}<\infty\) and \(X\) is
Banach, \(x:=\sum_r v_r\) exists in norm and \(x_R\to x\). For \(S>R\),
\(\abs{x_S-x_R}\le\sum_{r=R}^{S-1}\abs{v_r}\le T_R\); letting \(S\to\infty\),
Lemma~\ref{lem:norm_limit_abs_le} gives \(\abs{x-x_R}\le T_R\).

\emph{Step 3 (norm limits \(y_m\)).} Fix \(m\). For \(S>R\),
\(\abs{P_m x_S-P_m x_R}\le\sum_{r=R}^{S-1}\abs{P_m v_r}\le U_R\)
(Lemma~\ref{lem:proj_linear} and \(u_r\in A(v_r)\)); so \((P_m x_R)_R\) is norm
Cauchy. Let \(y_m=\lim_R P_m x_R\); then \(\abs{y_m-P_m x_R}\le U_R\) for all
\(m,R\) (Lemma~\ref{lem:norm_limit_abs_le}).

\emph{Step 4 (\(y_m\sotends x\)).} Fix \(R\). As \(P_m x_R\sotends x_R\)
(Lemma~\ref{lem:proj_sotends}) with some regulator \(w^{R}\dzero\), for each \(j\)
eventually \(\abs{P_m x_R-x_R}\le w^{R}_j\). Combining Steps 2--3,
\[
  \abs{y_m-x}\le\abs{y_m-P_m x_R}+\abs{P_m x_R-x_R}+\abs{x_R-x}
  \le U_R+T_R+w^{R}_j
\]
for all large \(m\). Set \(c_R=U_R+T_R\); then \(c\dzero\)
(Lemma~\ref{lem:reg_add}). Lemma~\ref{lem:diagonal} yields \(y_m\sotends x\).

\emph{Step 5 (\(y_m=P_m x\)).} Each \(P_m x_R\in E_m\), and \(E_m\) is closed
(Fact~\ref{fact:findim}), so \(y_m\in E_m\). The truncation \(Q_m\) is continuous
(finite-dimensional domain, Fact~\ref{fact:findim}) and
\(Q_m(P_{m+1}x_R)=P_m x_R\); passing to the limit, \(Q_m y_{m+1}=y_m\). With
Step 4, Lemma~\ref{lem:tower_eq_proj} gives \(y_m=P_m x\).

\emph{Step 6 (conclude).} By Steps 3 and 5,
\(\abs{P_m(x-x_R)}=\abs{P_m x-P_m x_R}=\abs{y_m-P_m x_R}\le U_R\) for all \(m\), so
\(U_R\in A(x-x_R)\) and \(\rho(x-x_R)\le\norm{U_R}\to0\) (Step 1). Hence
\(x_R\to x\) in \(\rho\).
\end{proof}

\begin{proposition}[\((X,\rho)\) is complete]\label{prop:rho_complete}
\leanname{instCompleteSpace\_Xrho}
\(\Xrho\) is a Banach space.
\uses{con:Xrho,lem:rho_series,fact:completecrit}
\end{proposition}
\begin{proof}
By Lemma~\ref{lem:rho_series} every \(\rho\)-norm-summable series converges in
\(\rho\); this is exactly the criterion of Fact~\ref{fact:completecrit}, so
\(\Xrho\) is complete.

\emph{Self-contained alternative} (matching the source note): a \(\rho\)-Cauchy
sequence has a subsequence \((x_r)\) with \(\sum_r\rho(x_{r+1}-x_r)<\infty\);
apply Lemma~\ref{lem:rho_series} to \(v_r=x_{r+1}-x_r\) to converge that
subsequence, and a Cauchy sequence with a convergent subsequence converges.
\end{proof}

%======================================================================
\section{The maximal inequality (Theorem A)}\label{sec:max}
%======================================================================

\begin{lemma}[Projections of a finite combination]\label{lem:proj_finite}
\leanname{proj\_of\_finite}
Fix \(N\ge1\), \(a:\N\to\R\), and put \(x=\sum_{k=0}^{N-1} a_k e_k\). Let
\(\hat a_k=a_k\) for \(k<N\) and \(\hat a_k=0\) otherwise. Then
\(a^{x}=\hat a\), and for all \(m\),
\[
  P_m x=\sum_{k=0}^{\min(m,N-1)} a_k e_k .
\]
In particular \(P_m x=s_m(a)\) for \(m\le N-1\) and \(P_m x=x\) for \(m\ge N-1\).
\uses{con:coord,lem:so_eventuallyEq,lem:coord_unique}
\end{lemma}
\begin{proof}
\(s_m(\hat a)=x\) for all \(m\ge N-1\), so \(s_m(\hat a)\sotends x\)
(Lemma~\ref{lem:so_eventuallyEq}); Lemma~\ref{lem:coord_unique} gives
\(a^{x}=\hat a\), and the formula for \(P_m x=s_m(\hat a)\) follows.
\end{proof}

\begin{lemma}[\(\rho\) of a finite combination is the maximal-function norm]\label{lem:rho_eq_sup}
\leanname{rho\_eq\_norm\_sup}
With \(N,a,x\) as in Lemma~\ref{lem:proj_finite} and
\[
  v=\bigjoin_{m=0}^{N-1}\Big|\sum_{k=0}^{m} a_k e_k\Big|
   \qquad\text{(Lean: \((\mathrm{range}\,N).\mathrm{sup}'\,(\lambda m.\,|\textstyle\sum_{k\in\mathrm{range}(m+1)}a_k\act e_k|)\)),}
\]
one has \(\rho(x)=\norm{v}\).
\uses{con:rho,lem:proj_finite,fact:finset,fact:solid,fact:csinf}
\end{lemma}
\begin{proof}
By Lemma~\ref{lem:proj_finite} the family \(\{\abs{P_m x}:m\in\N\}\) equals
\(\{\,|s_m(a)| : 0\le m\le N-1\,\}\) (values with \(m\ge N-1\) repeat
\(\abs{x}=|s_{N-1}(a)|\)). Thus \(v\) is its least upper bound
(Fact~\ref{fact:finset}): \(v\ge0\), \(\abs{P_m x}\le v\) for all \(m\)
(so \(v\in A(x)\)), and \(v\le u\) for every \(u\in A(x)\). Hence
\(\rho(x)\le\norm{v}\); and for \(u\in A(x)\), \(v\le u\) gives
\(\norm{v}\le\norm{u}\) (Fact~\ref{fact:solid}), so \(\norm{v}\le\rho(x)\)
(Fact~\ref{fact:csinf}). Therefore \(\rho(x)=\norm{v}\).
\end{proof}

\begin{lemma}[The comparison constant]\label{lem:const}
\leanname{exists\_const\_rho\_le}
There is \(C<\infty\) with \(\rho(x)\le C\norm{x}\) for all \(x\in X\).
\uses{con:Xrho,prop:rho_complete,lem:norm_le_rho,fact:banach}
\end{lemma}
\begin{proof}
The identity \(\iota:\Xrho\to X\) is linear and, by
Lemma~\ref{lem:norm_le_rho}, satisfies \(\norm{\iota x}=\norm{x}\le\rho(x)\), hence
is continuous with \(\norm{\iota}\le1\). It is bijective. Since \(\Xrho\) is
complete (Proposition~\ref{prop:rho_complete}) and \(X\) is Banach, the bounded
inverse theorem (Fact~\ref{fact:banach}, \texttt{ContinuousLinearEquiv.ofBijective}
with \(\ker\iota=\bot\), \(\operatorname{range}\iota=\top\)) makes
\(\iota^{-1}:X\to\Xrho\) continuous. With \(C=\norm{\iota^{-1}}\),
\(\rho(x)=\norm{\iota^{-1}x}_{\rho}\le C\norm{x}\).
\end{proof}

\begin{theorem}[Theorem A, \(\sigma\)-order form]\label{thm:main}
\leanname{theoremA\_sigma\_order\_form (FIXED name, from Showcase\_ThmA.lean)}
There is \(C<\infty\) such that for every \(N\ge1\) and all scalars
\(a_0,\dots,a_{N-1}\),
\[
  \Big\|\ \bigjoin_{m=0}^{N-1}\Big|\sum_{k=0}^{m} a_k e_k\Big|\ \Big\|
  \ \le\ C\,\Big\|\sum_{k=0}^{N-1} a_k e_k\Big\| .
\]
\uses{lem:const,lem:rho_eq_sup}
\end{theorem}
\begin{proof}
Take \(C\) from Lemma~\ref{lem:const}. Fix \(N,a\) and set
\(x=\sum_{k=0}^{N-1}a_k e_k\) and \(v\) as in Lemma~\ref{lem:rho_eq_sup}. Then
\(\norm{v}=\rho(x)\le C\norm{x}\), which is the claim.
(The spanning hypothesis \((\overline{\operatorname{span}}\,e=\top)\) of the Lean
statement is not used; see Remark~\ref{rem:hspan}.)
\end{proof}

\begin{remark}[On the unused spanning hypothesis]\label{rem:hspan}
The Lean statement carries \((\operatorname{span}_\R(\operatorname{range} e))\).
\texttt{topologicalClosure}\(=\top\). It plays no role in the proof: everything
follows from \texttt{IsSigmaOrderBasis} (Definition~\ref{def:basis}), which already
ranges over all \(x\in X\). Keep it in the signature (to match
\texttt{Showcase\_ThmA.lean}); it may raise an unused-variable warning only.
\end{remark}

%======================================================================
\section{Dependency graph (summary)}\label{sec:dag}
%======================================================================

Direct predecessors of each result (leaves \(=\) \S\ref{sec:mathlib}). Reading this
top to bottom is a valid topological order for formalization.

\begin{center}
\small
\begin{tabular}{ll}
\toprule
Result & Direct \texttt{uses} \\
\midrule
\ref{lem:smul_sup_nonneg} \texttt{smul\_sup\_of\_nonneg} & \ref{fact:possmul} \\
\ref{lem:abs_smul} \texttt{abs\_smul\_vectorLattice} & \ref{lem:smul_sup_nonneg}, \ref{fact:absgroup} \\
\ref{lem:reg_zero} \texttt{downToZero\_zero} & \ref{def:downto} \\
\ref{lem:reg_add} \texttt{DownToZero.add} & \ref{def:downto}, \ref{fact:absgroup} \\
\ref{lem:reg_smul} \texttt{DownToZero.smul} & \ref{def:downto}, \ref{lem:reg_zero}, \ref{fact:possmul} \\
\ref{lem:diagonal} \texttt{sigmaOrderTendsto\_of\_diagonal} & \ref{def:downto}, \ref{def:sigo}, \ref{lem:reg_add}, \ref{fact:absgroup}, \ref{fact:finset} \\
\ref{lem:so_eventuallyEq} & \ref{def:sigo}, \ref{lem:reg_zero} \\
\ref{lem:so_add} \texttt{SigmaOrderTendsto.add} & \ref{def:sigo}, \ref{lem:reg_add}, \ref{fact:absgroup} \\
\ref{lem:so_smul} \texttt{SigmaOrderTendsto.smul} & \ref{def:sigo}, \ref{lem:reg_smul}, \ref{lem:abs_smul}, \ref{fact:possmul} \\
\ref{lem:so_bounded} \texttt{.orderBounded} & \ref{def:sigo}, \ref{fact:absgroup}, \ref{fact:finset} \\
\ref{lem:so_abs_le} \texttt{.abs\_le} & \ref{def:sigo}, \ref{fact:absgroup} \\
\ref{lem:norm_limit_abs_le} \texttt{abs\_le\_of\_tendsto} & \ref{fact:orderclosed}, \ref{fact:absgroup} \\
\ref{lem:reg_tail} \texttt{downToZero\_tsum\_tail} & \ref{def:downto}, \ref{fact:series}, \ref{fact:solid} \\
\ref{con:coord} \texttt{coord, proj} & \ref{def:basis} \\
\ref{lem:proj_sotends} & \ref{con:coord}, \ref{def:basis} \\
\ref{lem:coord_unique} & \ref{con:coord}, \ref{def:basis} \\
\ref{lem:coord_add} / \ref{lem:coord_smul} & \ref{lem:proj_sotends}, \ref{lem:so_add}/\ref{lem:so_smul}, \ref{lem:coord_unique} \\
\ref{lem:proj_linear} & \ref{lem:coord_add}, \ref{lem:coord_smul} \\
\ref{lem:e_indep} & \ref{con:coord}, \ref{lem:so_eventuallyEq}, \ref{lem:coord_unique} \\
\ref{con:rho} \texttt{rho} & \ref{con:coord} \\
\ref{lem:rho_wd} & \ref{con:rho}, \ref{lem:proj_sotends}, \ref{lem:so_bounded}, \ref{fact:csinf} \\
\ref{lem:norm_le_rho} & \ref{con:rho}, \ref{lem:proj_sotends}, \ref{lem:so_abs_le}, \ref{fact:solid}, \ref{fact:csinf} \\
\ref{lem:rho_approx} & \ref{con:rho}, \ref{lem:rho_wd}, \ref{fact:csinf} \\
\ref{lem:rho_smul} & \ref{con:rho}, \ref{lem:proj_linear}, \ref{lem:abs_smul}, \ref{fact:possmul} \\
\ref{lem:rho_add} & \ref{con:rho}, \ref{lem:proj_linear}, \ref{lem:rho_approx}, \ref{fact:absgroup}, \ref{fact:csinf} \\
\ref{lem:rho_zero} & \ref{con:rho}, \ref{lem:norm_le_rho} \\
\ref{con:Xrho} \texttt{Xrho} & \ref{lem:rho_smul}, \ref{lem:rho_add}, \ref{lem:rho_zero}, \ref{fact:banach} \\
\ref{lem:tower_eq_proj} & \ref{lem:e_indep}, \ref{lem:coord_unique}, \ref{fact:findim} \\
\ref{lem:rho_series} & \ref{con:rho}, \ref{lem:rho_approx}, \ref{lem:norm_le_rho}, \ref{lem:reg_tail}, \ref{lem:reg_add}, \ref{lem:proj_sotends}, \ref{lem:proj_linear}, \ref{lem:norm_limit_abs_le}, \ref{lem:diagonal}, \ref{lem:tower_eq_proj}, \ref{fact:series}, \ref{fact:findim} \\
\ref{prop:rho_complete} & \ref{con:Xrho}, \ref{lem:rho_series}, \ref{fact:completecrit} \\
\ref{lem:proj_finite} & \ref{con:coord}, \ref{lem:so_eventuallyEq}, \ref{lem:coord_unique} \\
\ref{lem:rho_eq_sup} & \ref{con:rho}, \ref{lem:proj_finite}, \ref{fact:finset}, \ref{fact:solid}, \ref{fact:csinf} \\
\ref{lem:const} & \ref{con:Xrho}, \ref{prop:rho_complete}, \ref{lem:norm_le_rho}, \ref{fact:banach} \\
\ref{thm:main} \texttt{theoremA\_sigma\_order\_form} & \ref{lem:const}, \ref{lem:rho_eq_sup} \\
\bottomrule
\end{tabular}
\end{center}

%======================================================================
\section{Suggested split into Lean files}\label{sec:files}
%======================================================================

A layout mirroring the sections above. Import order is top-to-bottom; every file
imports \texttt{Mathlib} (or the relevant sub-imports) plus the files above it.
The three core definitions and the final theorem \emph{must} match
\texttt{Showcase\_ThmA.lean} verbatim; place them in \texttt{Defs.lean} and
\texttt{MaximalInequality.lean} respectively so the \texttt{Showcase} statement
type-checks unchanged against this development.

\begin{center}
\small
\begin{tabular}{lll}
\toprule
File & Contents (blocks) & Notes \\
\midrule
\texttt{ThmA/Prelim/AbsSmul.lean}
  & \ref{lem:smul_sup_nonneg}, \ref{lem:abs_smul}
  & the only new theory (Remark~\ref{rem:gap}) \\
\texttt{ThmA/Defs.lean}
  & \ref{def:downto}, \ref{def:sigo}, \ref{def:basis}, \ref{con:coord}
  & core defs \emph{verbatim} from Showcase \\
\texttt{ThmA/Regulators.lean}
  & \ref{lem:reg_zero}, \ref{lem:reg_add}, \ref{lem:reg_smul}, \ref{lem:diagonal}
  & needs \texttt{AbsSmul}? no; \texttt{Defs} only \\
\texttt{ThmA/SigmaOrder.lean}
  & \ref{lem:so_eventuallyEq}--\ref{lem:reg_tail}
  & Lemma 1 refactored; uses \texttt{AbsSmul}, \texttt{Regulators} \\
\texttt{ThmA/Coordinates.lean}
  & \ref{lem:proj_sotends}--\ref{lem:e_indep}
  & Lemma 2 refactored \\
\texttt{ThmA/AuxNorm.lean}
  & \ref{con:rho}--\ref{con:Xrho}
  & Lemma 3 + type synonym \(\Xrho\) \\
\texttt{ThmA/Complete.lean}
  & \ref{lem:tower_eq_proj}, \ref{lem:rho_series}, \ref{prop:rho_complete}
  & Proposition 4 (heaviest file) \\
\texttt{ThmA/MaximalInequality.lean}
  & \ref{lem:proj_finite}, \ref{lem:rho_eq_sup}, \ref{lem:const}, \ref{thm:main}
  & final theorem \emph{verbatim} from Showcase \\
\texttt{ThmA.lean}
  & (re-export) & imports all of the above \\
\bottomrule
\end{tabular}
\end{center}

\paragraph{Notes for the formalization agent.}
\begin{itemize}[leftmargin=1.5em]
\item Start with \texttt{AbsSmul.lean}: it unblocks \ref{lem:so_smul} and
      \ref{lem:rho_smul} and is the only place needing invention.
\item \texttt{Complete.lean} is the critical path; \ref{lem:rho_series} (six steps)
      and \ref{lem:tower_eq_proj} carry essentially all the difficulty. Formalize
      \ref{lem:tower_eq_proj} against Mathlib's \texttt{Basis}/\texttt{LinearIndependent}
      API first.
\item For \(\Xrho\) prefer building an \texttt{AddGroupNorm} bundle and calling
      \texttt{AddGroupNorm.toNormedAddCommGroup}, then supply
      \texttt{NormedSpace}\ \(\R\ \Xrho\) via \(\rho(\alpha\act x)=\abs{\alpha}\rho(x)\)
      (\ref{lem:rho_smul}); transport the \(\R\)-module structure through the type synonym.
\item \(\rho\) is best defined with \texttt{sInf} over \(\{\norm{u}\mid u\in A(x)\}\);
      keep \ref{lem:rho_wd} (nonempty + bdd-below) adjacent so \texttt{csInf} lemmas fire.
\item Do \emph{not} attempt \texttt{hspan} (Remark~\ref{rem:hspan}); it is inert.
\end{itemize}

\end{document}
